\section{ML}\label{sec:08-ml}

\begin{contentbox}
\textbf{1. Daten einlesen, Modell trainieren und validieren}

Dieser Code zeigt die Grundschritte eines typischen ML-Prozesses: Einlesen der Daten, Aufteilen in Trainings- und Testdaten, Trainieren eines Klassifikationsmodells (z.B. Decision Tree), sowie Validieren mittels Accuracy oder R2-Score.
\end{contentbox}
\begin{codebox}
\begin{lstlisting}[style=CodeExpert]
import pandas as pd
import numpy as np  # optional

# Daten einlesen und in X und y aufteilen
df = pd.read_csv("data.csv")
X = df.drop(["target"], axis=1)
y = df["target"]

# in Trainings- und Testdaten aufteilen
from sklearn.model_selection import train_test_split
X_train, X_test, y_train, y_test =
train_test_split(X, y, test_size=0.2)

# Modell trainieren
from sklearn.tree import DecisionTreeClassifier
model = DecisionTreeClassifier()
model.fit(X_train, y_train)

# Klassifikationsmodell validieren
from sklearn.metrics import accuracy_score
y_pred = model.predict(X_test)
validation_score = accuracy_score(y_test, y_pred)

# oder Regressionsmodell validieren
from sklearn.metrics import r2_score
validation_score = r2_score(y_test, y_pred)
\end{lstlisting}
\end{codebox}

\begin{contentbox}
\textbf{2. Vorhersage mit neuen Daten}

Sobald ein Modell trainiert ist, kann es mit neuen Eingabedaten (\texttt{X\_final}) verwendet werden, um Vorhersagen (\texttt{y\_final}) zu erzeugen.
\end{contentbox}
\begin{codebox}
\begin{lstlisting}[style=CodeExpert]
X_final = pd.read_csv("X_final.csv")
y_final = model1.predict(X_final)
return y_final
\end{lstlisting}
\end{codebox}

\begin{contentbox}
\textbf{3. Neuronales Netz mit \texttt{MLPClassifier}}

Ein Beispiel für ein neuronales Netz (Multilayer Perceptron). Die Architektur wird über \texttt{hidden\_layer\_sizes} definiert, z.B. fünf versteckte Schichten mit je 6–8 Neuronen.
\end{contentbox}
\begin{codebox}
\begin{lstlisting}[style=CodeExpert]
from sklearn.neural_network import MLPClassifier

X_train, X_test, y_train, y_test = train_test_split(
    X, y, test_size=0.1, random_state=31)

model = MLPClassifier(
    hidden_layer_sizes=(8, 6, 7, 6, 7, 5),
    max_iter=5000,
    random_state=85)

model.fit(X_train, y_train)
\end{lstlisting}
\end{codebox}

\subsection{Wichtige Parameter}\label{sec:08-wichtige-parameter}
\begin{contentbox}
\begin{itemize}
  \item \textbf{max\_iter}: Maximale Anzahl Iterationen, die ein Modell während des Trainings durchläuft, um die Verlustfunktion zu minimieren (z.\,B. Training endet, wenn Loss $< 0.001$).\\
  \textit{Verwendet von:} \texttt{LogisticRegression}, \texttt{MLPClassifier}, \texttt{MLPRegressor}

  \item \textbf{hidden\_layer\_sizes}: Gibt die Struktur eines neuronalen Netzes an (z.\,B. Anzahl Neuronen pro Schicht).\\
  \textit{Verwendet von:} \texttt{MLPClassifier}, \texttt{MLPRegressor}

  \item \textbf{random\_state}: Initialisiert den Zufallsprozess (z.\,B. Shuffle bei \texttt{train\_test\_split}, Gewichtung im Netz), um reproduzierbare Ergebnisse zu garantieren.\\
  \textit{Verwendet von:} allen Modulen mit Zufallskomponenten

  \item \textbf{test\_size}: Prozentualer Anteil der Daten, der für die Testmenge reserviert wird. Üblich ist $0.1$ bis $0.3$.\\
  \textit{Verwendet von:} \texttt{train\_test\_split}
 \item \textbf{cv} (Cross-Validation): Gibt bei \texttt{GridSearchCV} oder \texttt{cross\_val\_score} die Anzahl an Folds an, mit denen K-Fold-Cross-Validation durchgeführt wird.\\
  \textit{Verwendet von:} \texttt{GridSearchCV}, \texttt{cross\_val\_score}, etc.

  \item \textbf{n\_clusters}: Gibt an, wie viele Cluster ein Clustering-Algorithmus bilden soll. Wichtig bei unüberwachtem Lernen.\\
  \textit{Verwendet von:} \texttt{KMeans}, \texttt{DBSCAN (implizit)}, \texttt{AgglomerativeClustering}

  \item \textbf{Grid Search via Cross-Validation}:
	\textit{Normales Grid Search:}
Du probierst verschiedene Modelltiefen und weitere Hyperparameter-Kombinationen aus.
Für jede Kombination misst du die Fehlerrate des Modells.
	\textit{k-Fold-Cross-Validation:}
Die Daten werden in $k$ gleich grosse Teile (Folds) aufgeteilt.
Das Modell wird $k$-mal trainiert — jedes Mal mit $k-1$ Folds als Trainingsdaten und dem verbleibenden Fold als Validierungsmenge.
Jeder Fold dient genau einmal als Validierung, sodass eine zuverlässige Schätzung der Modellgüte entsteht.
\end{itemize}
\end{contentbox}

\subsection{Modelle}\label{sec:08-modelle}

\subsubsection{ML-Typen: Regression vs. Klassifikation}
\begin{contentbox}
\textbf{Regression} wird verwendet, um stetige Werte vorherzusagen (z.\,B. Preis, Wahrscheinlichkeit) in '\%'. \\
\textbf{Klassifikation} hingegen ordnet Datenpunkte einer festen Klasse zu (z.\,B. Spam/Nicht-Spam), ja oder nein.
\end{contentbox}



\subsubsection{Modelle für Regression}

\begin{contentbox}[Linear Regression]
\begin{itemize}
  \item Vorhersage basiert auf einer linearen Kombination der Eingabefeatures:
  \[
  y = w_1 x_1 + w_2 x_2 + \cdots + w_n x_n + b
  \]
  \item \textbf{Import:} \texttt{from sklearn.linear\_model import LinearRegression}
\end{itemize}
\end{contentbox}

\begin{contentbox}[Decision Tree Regressor]
\begin{itemize}
  \item Modelliert nicht-lineare Beziehungen durch Entscheidungen an Verzweigungsknoten.
  \item Zerlegt die Eingabemenge in Zonen mit möglichst homogener Zielvariable.
  \item \textbf{Import:} \texttt{from sklearn.tree import DecisionTreeRegressor}
\end{itemize}
\end{contentbox}

\begin{contentbox}[Neural Network Regressor]
\begin{itemize}
  \item Ein neuronales Netz mit Regressionsziel: nicht-linear, viele Parameter.
  \item Flexible Anpassung durch \texttt{hidden\_layer\_sizes}, \texttt{max\_iter}, \texttt{activation}.
  \item \textbf{Import:} \texttt{from sklearn.neural\_network import MLPRegressor}
\end{itemize}
\end{contentbox}

\subsubsection{Bewertung mit dem R2-Score}
\begin{contentbox}
Der $R^2$-Score misst die Qualität einer Regressionsvorhersage:

\[
R^2(D,f) = 1 - \frac{MSE(D,f)}{MSE(D,f_0)}
\]

\begin{itemize}
  \item \textbf{MSE hängt von der Skalierung der Daten ab, $R^2$ - Score nicht}
  \item $f_0$: konstante Durchschnittsvorhersage
  \item $D$: Datensatz
  \item $MSE$: Mean Squared Error, misst mittlere quadratische Abweichung (Verlustfunktion):
  \[
  MSE(D,f) = \frac{1}{n} \sum_{i=1}^n (y_i - \hat{y}_i)^2,\quad \hat{y}_i = f(x_i)
  \]
\end{itemize}
\end{contentbox}


\subsubsection{LinearRegression}
\begin{contentbox}
    \textbf{Wann einsetzen}
    Für einfache Regressionsprobleme, wenn der Zusammenhang zwischen Features und Zielwert annähernd linear ist.

\textbf{Stärken und Schwächen}
\textbf{Stärken:} schnell, gut interpretierbar, theoretisch gut verstanden.
\textbf{Schwächen:} empfindlich gegenüber Ausreissern, setzt Linearität voraus, kann komplexe Muster nicht abbilden.

\textbf{Wichtige Hyperparameter}
\texttt{fit\_intercept}: ob ein Bias-Term berücksichtigt wird.
\texttt{n\_jobs}: Anzahl CPU-Kerne, die während des Fits genutzt werden.

\textbf{Beispiel-Use-Case}
Vorhersage von Hauspreisen anhand von Features wie Wohnfläche, Anzahl Zimmer und Baujahr in einer kleinen Nachbarschaft.
\end{contentbox}
\begin{codebox}
\begin{lstlisting}[style=CodeExpert]
from sklearn.linear_model import LinearRegression

model = LinearRegression()
model.fit(X_train, y_train)
y_pred = model.predict(X_test)
\end{lstlisting}
\end{codebox}

% -----------------------------


% -----------------------------
\subsubsection{Modelle für Klassifikation}

\begin{statementbox}[Voraussetzung]
Die Daten müssen gelabelt sein.
\end{statementbox}

\begin{contentbox}[Logistic Regression]
\begin{itemize}
  \item Verwendet die \textbf{sigmoid}-Funktion zur Klassenvorhersage:
  \[
  y = \sigma(w_0 + w_1 x_1 + \cdots + w_n x_n)
  \]
  \item Ausgabe ist die Wahrscheinlichkeit, dass ein Punkt zur positiven Klasse gehört.
  \item \textbf{Import:} \texttt{from sklearn.linear\_model import LogisticRegression}
\end{itemize}
\end{contentbox}

\begin{contentbox}[Decision Tree Classifier]
\begin{itemize}
  \item Genau wie der Regressor, aber mit Klassenzuordnung am Ende der Blätter.
  \item Tiefe entspricht max. Anzahl der Entscheidungen
  \item Gut interpretierbar, aber anfällig für Overfitting.
  \item \textbf{Import:} \texttt{from sklearn.tree import DecisionTreeClassifier}
\end{itemize}
\end{contentbox}

\begin{contentbox}[Neural Network Classifier]
\begin{itemize}
  \item Tieferes neuronales Netz, spezialisiert auf Klassifikationsaufgaben.
  \item Meist verwendet: \texttt{MLPClassifier} (Multilayer Perceptron)
  \item \textbf{Import:} \texttt{from sklearn.neural\_network import MLPClassifier}
\end{itemize}
\end{contentbox}



\subsubsection{(Balanced) Accuracy}
\begin{contentbox}


\textbf{Confusion Matrix:}

\begin{center}
\begin{ZSFtablePlain}{C C C}
\toprule
\textbf{real / vorhergesagt} & \textbf{Positiv} & \textbf{Negativ} \\
\midrule
\textbf{Positiv} & True Pos. (TP) & False Neg. (FN) \\
\textbf{Negativ} & False Pos. (FP) & True Neg. (TN) \\
\bottomrule
\end{ZSFtablePlain}
\end{center}



\textbf{Accuracy:} Anteil der korrekt klassifizierten Beobachtungen:

\begin{equation}
\text{Accuracy} = \frac{TP + TN}{TP + TN + FP + FN}.
\end{equation}

\textbf{Balanced Accuracy:} Mittelwert der True Positive Rate (TPR) und der True Negative Rate (TNR): \\
Aka Zeilen werden Normiert (Summe jeder Zeile ist 1)

\begin{equation}
\text{Balanced Accuracy} = \frac{1}{2} \big( \text{TPR} + \text{TNR} \big)
\end{equation}

mit:

\begin{equation}
\text{TPR} = \frac{TP}{TP + FN}
\quad \text{und} \quad
\text{TNR} = \frac{TN}{TN + FP}.
\end{equation}

\ZSFgapL

\textit{Hinweis:} Die Balanced Accuracy ist nützlich bei unausgeglichenen Klassenverteilungen, da beide Klassen gleich stark berücksichtigt werden.
\end{contentbox}



\subsubsection{0-1-Loss:}
\begin{contentbox}


Die \textit{0-1-Loss-Funktion} misst, ob eine Vorhersage korrekt ist:

\begin{equation}
L(y, \hat{y}) =
\begin{cases}
0, & \text{wenn } y = \hat{y} \\
1, & \text{wenn } y \neq \hat{y}.
\end{cases}
\end{equation}

Für einen Datensatz mit $n$ Beispielen:

\begin{equation}
\text{Empirical Risk} = \frac{1}{n} \sum_{i=1}^{n} L(y_i, \hat{y}_i).
\end{equation}
Im Bspl der Confusion Matrix:
\begin{equation}
\text{0-1-Loss (Error Rate)} = \frac{FP + FN}{TP + TN + FP + FN}.
\end{equation}
Die Accuracy ist das Komlement zum Loss
\begin{equation}
    \text{Accuracy} = 1 - \text{Loss} = \frac{TP + TN}{TP + TN + FP + FN}.
\end{equation}
\end{contentbox}

\subsubsection{Gini-index Loss:}
\begin{contentbox}


\textbf{Gini-Index (als Loss im Entscheidungsbaum):}

Die „Impurity“ einer Partition $Part_m$ ist:
\end{contentbox}
\begin{contentbox}
    

\begin{equation}
G(Part_m) = \sum_{k} \hat{p}_{m,k} (1 - \hat{p}_{m,k}),
\end{equation}

wobei $\hat{p}_{m,k}$ der Anteil der Klasse/Labels $k$ in der Partition ist.



Ein \textbf{Split} teilt den Knoten $Part_m$ in z.B. zwei Partitionen $Part_1$ und $Part_2$.

Der Gini-Index für den gesamten Split ist:

\begin{equation}
G_{\text{split}} =
\frac{|Part_1|}{|Part_m|} G(Part_1)
+ \frac{|Part_2|}{|Part_m|} G(Part_2).
\end{equation}

wobei $|Part_i|$ die Anzahl der Datenpunkte in der Partition $i$ ist

\ZSFgapL

\textit{Hinweis:} Ein Entscheidungsbaum wählt den Split, der den Gini-Index minimiert.

\end{contentbox}
\subsubsection{LogisticRegression}
\begin{contentbox}
\textbf{Wann einsetzen}
Für binäre oder multinomiale Klassifikation, wenn die Entscheidungsgrenze annähernd linear ist. Liefert Wahrscheinlichkeitsschätzungen.

\textbf{Stärken und Schwächen}
\textbf{Stärken:} interpretierbar, schnell, liefert Klassenwahrscheinlichkeiten.
\textbf{Schwächen:} schwierig bei nicht-linearen Mustern.

\textbf{Wichtige Hyperparameter}

\texttt{max\_iter}: ca. 500--1000, \texttt{class\_weight}: \texttt{"balanced"}, \texttt{None}.

\textbf{Beispiel-Use-Case}
Vorhersage, ob ein Patient Diabetes hat (ja/nein), basierend auf medizinischen Messwerten — wenn die Features einen linearen Zusammenhang aufweisen.
\end{contentbox}
\begin{codebox}
\begin{lstlisting}[style=CodeExpert]
from sklearn.linear_model import LogisticRegression

model = LogisticRegression(max_iter = 1000,
                        "class_weight" = balanced)
model.fit(X_train, y_train)
y_pred = model.predict(X_test)
\end{lstlisting}
\end{codebox}

\subsubsection{DecisionTreeClassifier}
\begin{contentbox}
\textbf{Faustregeln}
\texttt{max\_depth}: 3--10 für kleine/mittlere Datensätze. Beschränken, um Overfitting zu vermeiden.
\texttt{min\_samples\_split}: erhöhen (z.\,B. 10+), um zu regularisieren.
\texttt{min\_samples\_leaf}: 1--5, um Vorhersagen zu glätten.
\end{contentbox}
\begin{codebox}
\begin{lstlisting}[style=CodeExpert]
from sklearn.tree import DecisionTreeClassifier

model = DecisionTreeClassifier(
    # aus Vorlesung
    # für Reproduzierbarkeit
    random_state=0
    # Modellkomplexität steuern
    max_depth=4,
    # optional
    # "entropy" als Alternative
    criterion="gini",
    splitter="best",          # oder "random"
    min_samples_split=2,      # Standard
    min_samples_leaf=1,       # Standard
    # "sqrt" oder "log2" probieren
    max_features=None,

)
model.fit(X_train, y_train)
\end{lstlisting}
\end{codebox}

% -----------------------------

\subsubsection{MLPClassifier}
\begin{contentbox}
\textbf{Faustregeln}
\texttt{hidden\_layer\_sizes}: 1--2 Schichten mit je 50--100 Neuronen funktionieren oft gut.
\texttt{activation}: \texttt{"relu"} verwenden. \texttt{"tanh"} bei kleinen/strukturierten Daten.
\texttt{alpha}: Default ist 0.0001. Erhöhen, um Overfitting zu reduzieren.
\texttt{max\_iter}: hoch setzen (z.\,B. 500), falls das Training nicht konvergiert.
\end{contentbox}

\begin{codebox}
\begin{lstlisting}[style=CodeExpert]
from sklearn.neural_network import MLPClassifier

model = MLPClassifier(
    # zwei Schichten
    hidden_layer_sizes=(50, 50),
    # erhöhen falls nötig
    max_iter=5000,
    random_state=0
    # optional
    # oder "tanh", "logistic"
    activation="relu",
    # oder "sgd", "lbfgs"
    solver="adam",
    # L2-Regularisierung
    alpha=0.0001,
    # oder "adaptive"
    learning_rate="constant",
)
model.fit(X_train, y_train)
\end{lstlisting}
\end{codebox}

% -----------------------------

\subsubsection{GridSearchCV}
\begin{contentbox}
\textbf{Faustregeln}
\texttt{param\_grid}: 2--4 Werte pro Parameter.
\texttt{cv}: 5-Fold ist Standard. 3 für Geschwindigkeit, 10 für höhere Genauigkeit.
\texttt{scoring}: Metrik passend zur Aufgabe wählen.
\end{contentbox}
\begin{codebox}
\begin{lstlisting}[style=CodeExpert]
from sklearn.model_selection import GridSearchCV
from sklearn.tree import DecisionTreeClassifier

model = DecisionTreeClassifier(random_state=0)

param_grid = {
    "max_depth": [2, 4, 6],
    "class_weight": [None, "balanced"]
    "min_samples_split": [2, 10],
    "min_samples_leaf": [1, 5]

} """param_grid kann auch für weitere Parameter
wie alpha genutzt werden (Prüfung 02.25)"""


search = GridSearchCV(
    estimator=model,
    param_grid=param_grid,
    # oder "accuracy", "recall", etc.
    scoring="balanced_accuracy",
    cv=5,
    verbose=1    # nur Terminalausgabe
)
search.fit(X_train, y_train)
\end{lstlisting}
\end{codebox}



\subsection{Nicht-numerische Daten: Kategorische Merkmale encodieren}\label{sec:08-kategorische-daten}
\begin{contentbox}
Maschinelles Lernen erfordert in der Regel numerische Eingabedaten. Nicht-numerische (kategorische) Daten wie z.\,B. \texttt{“Farbe”} oder \texttt{“Beruf”} müssen zuerst in eine numerische Form gebracht werden. Dazu gibt es verschiedene Kodierungsverfahren – je nach Kontext und Modell mit eigenen Vor- und Nachteilen:
\end{contentbox}

\begin{contentbox}[Ordinal Encoding]
\textbf{Pro:}
\begin{itemize}
  \item Einfach umsetzbar.
  \item Ermöglicht die Abbildung von geordneter Struktur (z.\,B. 1 = Crew, 2 = First Class).
\end{itemize}
\textbf{Con:}
\begin{itemize}
  \item Kann eine falsche Nähe zwischen Kategorien suggerieren (z.\,B. 2 ist näher an 1 als 3, obwohl inhaltlich 1 und 3 ähnlicher sind).
  \item Führt möglicherweise zu verzerrten Modellinterpretationen.
\end{itemize}
\end{contentbox}

\begin{contentbox}[Mean Encoding]
\textbf{Pro:}
\begin{itemize}
  \item Nutzt den Mittelwert der Zielvariable pro Kategorie – kann nützliche Informationen übertragen.
  \item Versucht „scheinbare Nähe” zwischen Kategorien auszugleichen.
\end{itemize}
\textbf{Con:}
\begin{itemize}
  \item Gefahr von \ZSFkeyword{Data Leakage}: Die Zielvariable „sickert durch” in die Features.
  \item Erhöhtes Risiko von Overfitting.
  \item Verlangt Rechenaufwand.
\end{itemize}
\end{contentbox}

\begin{contentbox}[One-hot Encoding]
\textbf{Pro:}
\begin{itemize}
  \item Höchste Präzision: Kodiert jede Kategorie in einer separaten Spalte.
  \item Besonders geeignet für nicht-ordinale Kategorien.
\end{itemize}
\textbf{Con:}
\begin{itemize}
  \item Jeder neue Wert erzeugt eine neue Spalte (hoher Speicherplatzverbrauch).
  \item Bei vielen Kategorien kann das Datenmodell sehr breit werden („Fluch der Dimensionalität”).
\end{itemize}
\end{contentbox}
\begin{codebox}
\begin{lstlisting}[style=CodeExpert]
import pandas as pd

pd.get_dummies(a)
\end{lstlisting}
\end{codebox}


\subsection{Dimensionsreduktion}\label{sec:08-dimensionsreduktion}
\begin{contentbox}
\textbf{Faustregeln}
\texttt{n\_components}: 2--3 für Visualisierung, höher für Preprocessing.
\texttt{PCA}: linear, schnell, maximiert Varianz.
\texttt{t-SNE}: nichtlinear, gut für Cluster-Visualisierung, langsamer.
\end{contentbox}
\begin{codebox}
\begin{lstlisting}[style=CodeExpert]
from sklearn.decomposition import PCA
from sklearn.manifold import TSNE
from sklearn.preprocessing import StandardScaler

# --- Vorverarbeitung ---
# Normalisieren (bei Bilddaten üblich)
images_normalized = images / 255.0

# Oder Standardisieren (bei tabellarischen Daten üblich)
scaler = StandardScaler()
X_scaled = scaler.fit_transform(X)

# Bilder "flatten" (z.B. 28x28 -> 784 Dimensionen)
images_flat = images.reshape(len(images), -1)

# --- Dimensionsreduktion ---
# PCA (linear)
model = PCA(n_components=2)

# oder t-SNE (nichtlinear, gleiche API)
# model = TSNE(n_components=2, random_state=0)

# Fit und Transformation
images_trans = model.fit_transform(images_flat)
\end{lstlisting}
\end{codebox}
\subsection{Clustering}\label{sec:08-clustering}

\begin{contentbox}
\textbf{Faustregeln}
\texttt{StandardScaler}: immer vor PCA anwenden, wenn Features unterschiedliche Skalen haben.
\texttt{n\_components}: 2--3 für Visualisierung, mehr für Analyse.
\texttt{K-Means}: \texttt{init="k-means++"} liefert stabilere Ergebnisse.
\end{contentbox}
\begin{codebox}
\begin{lstlisting}[style=CodeExpert]
import numpy as np
from sklearn.datasets import load_digits
from sklearn.cluster import KMeans
from sklearn.preprocessing import StandardScaler
from sklearn.decomposition import PCA

# --- Daten laden ---
data, labels = load_digits(return_X_y=True)
(n_samples, n_features), n_digits = data.shape, np.unique(labels).size

# --- Vorverarbeitung ---
scaled = StandardScaler().fit_transform(data)

# --- PCA-Reduktion auf 2 Dimensionen ---
reduced_data = PCA(n_components=2).fit_transform(scaled)

# --- K-Means Clustering ---
kmeans = KMeans(init="k-means++", n_clusters=n_digits, n_init=4)
kmeans.fit(reduced_data)
\end{lstlisting}
\end{codebox}



\subsection{Was tun, wenn der \texttt{accuracy\_score} zu niedrig ist?}\label{sec:08-accuracy-zu-niedrig}

\begin{contentbox}
Wenn dein Klassifikationsmodell eine zu geringe Genauigkeit erreicht, kannst du systematisch folgende Maßnahmen ausprobieren:

\begin{enumerate}
  \item \textbf{Mehr Daten:} 
  \begin{itemize}
    \item Größere Trainingsmengen helfen oft, Muster robuster zu erkennen.
    \item Datenaugmentation oder neue Datenquellen prüfen.
  \end{itemize}

  \item \textbf{Feature Engineering:}
  \begin{itemize}
    \item Relevante Merkmale hinzufügen, irrelevante entfernen.
    \item Skalierung oder Transformation (z.\,B. \texttt{StandardScaler}).
    \item Kategorische Variablen korrekt encodieren (\texttt{One-Hot}, \texttt{Ordinal}, \texttt{Mean}).
  \end{itemize}

  \item \textbf{Modellparameter optimieren:}
  \begin{itemize}
    \item \texttt{hidden\_layer\_sizes} beim MLPClassifier vergrößern oder anpassen.
    \item \texttt{max\_iter} erhöhen (z.\,B. 2000 statt 500), damit das Modell besser konvergiert.
    \item \texttt{alpha} bei Regularisierung (MLP, Ridge\ldots) anpassen.
    \item Andere \texttt{activation} Funktionen testen: \texttt{"tanh"}, \texttt{"relu"}\ldots
  \end{itemize}

  \item \textbf{Randomisierung kontrollieren:}
  \begin{itemize}
    \item \texttt{random\_state} anpassen, wenn das Ergebnis stark schwankt.
    \item Mehrere Splits testen und Mittelwert bilden (z.\,B. mit \texttt{cross\_val\_score}).
  \end{itemize}

  \item \textbf{Anderes Modell testen:}
  \begin{itemize}
    \item Statt Decision Tree → \texttt{RandomForestClassifier} oder \texttt{MLPClassifier}.
    \item Für lineare Probleme eventuell \texttt{LogisticRegression}.
  \end{itemize}

  \item \textbf{Overfitting prüfen:}
  \begin{itemize}
    \item Wenn Training sehr gut, aber Test schlecht: Modell zu komplex?
    \item \texttt{max\_depth} bei Trees begrenzen, \texttt{dropout} oder \texttt{alpha} bei MLP verwenden.
  \end{itemize}
\end{enumerate}
\end{contentbox}




